\chapter{Attention Mechanisms}

The attention mechanism is the only component that operates over all
token pairs simultaneously. Chapters 4 and 5 asked how to enrich the
information arriving at the attention layer; this chapter asks whether
the attention computation itself should be modified. The central
question is whether differential attention,which computes the
difference of two softmax maps to suppress diffuse, low-information
attention patterns, transfers from its original NLP setting to
particle classification.

The physics motivation is direct: in a multi-top event, only a small
subset of the 10--20 reconstructed objects are decay products of the
four top quarks. Standard softmax attention distributes probability
mass across all token pairs including unrelated soft objects.
Differential attention is designed for exactly this setting, and
provides a clear testable hypothesis.

Two experiments are run. Exp.\ 6A establishes the headline comparison
across attention type, internal normalisation, and the presence of a
Lorentz-scalar bias. Exp.\ 6B is a targeted follow-up on how physics
biases interact with the two differential attention branches.

% ──────────────────────────────────────────────
\section{Attention type and internal normalisation (Exp.\ 6A)}
\textit{Lens: performance and efficiency.}

Sweeps A3 (standard vs.\ differential), A4 (attention-internal
normalisation), and whether a Lorentz-scalar bias is active. The bias
inclusion tests whether physics structure in the attention logits and
noise suppression from differential attention are complementary or
redundant. Wall-clock time per epoch is recorded alongside AUROC, since
differential attention runs two softmax passes per head.

\begin{table}[h]
\centering\small
\begin{tabular}{llr}
\toprule
Axis & Values & Config \\
\midrule
A3 — Attention type & \texttt{standard}, \texttt{differential} & 2 \\
A4 — Internal norm  & \texttt{none}, \texttt{layernorm}, \texttt{rmsnorm} & 3 \\
B1 — Lorentz bias   & \texttt{none}, \texttt{lorentz\_scalar} & 2 \\
Seeds & & 3 \\
\midrule
Total & & \textbf{36} \\
\bottomrule
\end{tabular}
\caption{Exp.\ 6A. All other axes at chapter baseline.
\texttt{lorentz\_scalar} uses the best feature set from Ch.\,5.}
\end{table}

Outputs: seed-averaged AUROC bar chart for all 12 configurations, with
wall-clock time as a secondary axis. Attention maps extracted at
inference time for standard and differential (both with A4=none,
B1=none), averaged over a fixed held-out batch and labelled by particle
type. If differential attention is working as intended, maps should show
higher concentration on physically proximate pairs.

% ──────────────────────────────────────────────
\section{Differential bias mode (Exp.\ 6B)}
\textit{Lens: performance.}

Restricted to A3=\texttt{differential}. Tests how the additive
Lorentz-scalar bias from Ch.\,5 should interact with the two attention
branches: added to neither, shared across both, or split with separate
parameters per branch.

\begin{table}[h]
\centering\small
\begin{tabular}{llr}
\toprule
Axis & Values & Config \\
\midrule
A3-a — Diff bias mode & \texttt{none}, \texttt{shared}, \texttt{split} & 3 \\
Seeds & & 3 \\
\midrule
Total & & \textbf{9} \\
\bottomrule
\end{tabular}
\caption{Exp.\ 6B. A3=\texttt{differential} and
B1=\texttt{lorentz\_scalar} fixed throughout.}
\end{table}

Output: compact three-row AUROC table. If \texttt{shared} and
\texttt{split} perform equivalently, the second branch does not require
independent bias information and the simpler shared parameterisation
is sufficient.

\subsection*{Chapter summary}

\begin{table}[h]
\centering\small
\begin{tabular}{llr}
\toprule
Experiment & Primary axes & Runs \\
\midrule
6A — Attention type and norm & A3, A4, B1 & 36 \\
6B — Differential bias mode  & A3-a        &  9 \\
\midrule
\textbf{Total} & & \textbf{45} \\
\bottomrule
\end{tabular}
\end{table}
